\documentclass[
	ngerman,
    %draft, %Dient zur Erkennung von 'overfull-Boxen' + Verschnellert das Kompilierens, da Bilder nicht aufgelöst werden + Keine Links
	toc=listof, % Abbildungsverzeichnis sowie Tabellenverzeichnis in das Inhaltsverzeichnis aufnehmen
	toc=bibliography, % Literaturverzeichnis in das Inhaltsverzeichnis aufnehmen
	footnotes=multiple, % Trennen von direkt aufeinander folgenden Fußnoten
	parskip=half, % vertikalen Abstand zwischen Absätzen verwenden anstatt horizontale Einrückung von Folgeabsätzen
	numbers=noendperiod, % Den letzten Punkt nach einer Nummerierung entfernen (nach DIN 5008)
]{scrartcl}


%Metadaten der Dokumentation

%% Projektdaten ------------------------------------------------------
%Projekttitel 
\newcommand{\projektKurztitel}{Kurztitel\xspace}
\newcommand{\projektLangtitel}{Langtitel\xspace}


%Hausarbeit etc.
\newcommand{\kategorie}{KategorieDerArbeit\xspace}

\newcommand{\startTerminProjekt}{24.06.2025}
\newcommand{\abgabetermin}{01.01.9999}
\newcommand{\pruefungsID}{XXXXX\_P}


%%Autoren ------------------------------------------------------
%Autor One definieren
\newcommand{\autorVornameOne}{Max\xspace}
\newcommand{\autorNachnameOne}{Mustermann\xspace}
\newcommand{\autorFullNameOne}{\autorVornameOne \autorNachnameOne\xspace}
\newcommand{\autorFullNameOneOhneLeerzeichen}{\autorVornameOne \autorNachnameOne}
\newcommand{\autorAnschriftOne}{Musterstraße 13\xspace}
\newcommand{\autorOrtOne}{12345 Musterstadt\xspace}


%Weitere Daten Autor One
\newcommand{\autorMatrikelnummerOne}{12345\xspace}
\newcommand{\autorSemesterNummerOne}{-1\xspace}
\newcommand{\autorHandyPrivatOne}{123456789012\xspace}
\newcommand{\autorMailPrivatOne}{max-mustermann@domain.de\xspace}
\newcommand{\autorMailWorkOne}{muster@mail.de\xspace}

%Weiteres
\newcommand{\autorenFullNameRow}{\autorFullNameOne}



%Ansprechpartner / Dozenten ------------------------------------------------------
\newcommand{\ansprechpartnerOne}{Prof. Dr.-Ing. Max Mustermann\xspace}
\newcommand{\ansprechpartnerTwo}{Erika Mustermann \xspace}

\newcommand{\ansprechpartnerRow}{\ansprechpartnerOne, \ansprechpartnerTwo}



%%Hochschule ------------------------------------------------------
\newcommand{\studiengang}{NameStudiengang\xspace}
\newcommand{\vorlesungTitel}{VorlesungTitel\xspace}


%Hochschule Logos
\newcommand{\betriebHSHarzShort}{02_Datein/01_Bilder/00_Base/LogoHSHarz_Short.png}
\newcommand{\betriebHSHarzFull}{02_Datein/01_Bilder/00_Base/LogoHochschule_Full.png}


\newcommand{\eidesstattlicheErklaerungDatum}{\today}
\newcommand{\eingereichtAmDatum}{\today} % Metadaten zu diesem Dokument (Autor usw.)
%Für Sonderzeichen und Umlaute
\usepackage[T1]{fontenc}
\usepackage[utf8]{inputenc}

%Schrift ändern
\usepackage[scaled]{uarial} %Klon von Arial - mit Serifen

\renewcommand\familydefault{\sfdefault} %Default-Schrift auf "ohne Serifen" umstellen

% Schriftgröße zulassen
\usepackage{anyfontsize}

% Dokument auf Deutsch umstellen
\usepackage{babel}

% Tabellen
\PassOptionsToPackage{table}{xcolor}
\usepackage{tabularx}
\usepackage{colortbl}
\usepackage{float} %Damit kann man den Parameter "H" bei der table Umgebung verwenden und die Tabelle bleibt an der Stelle.

% Tabellen Linien formatieren
\usepackage{tabularx}
\usepackage{bigstrut}

% Einfache Definition der Zeilenabstände und Seitenränder etc.
\usepackage{setspace}
\usepackage{geometry}
%\usepackage[showframe]{geometry} %Zeigt einen Rahmen um den Text. Für Überlappungen.

% Einstellungen für PDF-Dateien
\usepackage{pdfpages}
\usepackage[
    bookmarks,
    bookmarksnumbered,
    bookmarksopen=true,
    bookmarksopenlevel=1,
    colorlinks=true,
    % di%ese Farbdefinitionen zeichnen Links im PDF farblich aus
    anchorcolor=MainColor,% Ankertext
    citecolor=MainColor, % Verweise auf Literaturverzeichniseinträge im Text
    filecolor=MainColor, % Verknüpfungen, die lokale Dateien öffnen
    menucolor=MainColor, % Acrobat-Menüpunkte
    urlcolor=MainColor,
    linkcolor=MainColor,
    %%
    pdftex,
    plainpages=false, % zur korrekten Erstellung der Bookmarks
    pdfpagelabels=true, % zur korrekten Erstellung der Bookmarks
    hypertexnames=false, % zur korrekten Erstellung der Bookmarks
    linktoc=page,
]{hyperref}
\hypersetup{
    pdftitle={\projektKurztitel -- \projektLangtitel},
    pdfauthor={\autorFullNameOne},
    pdfcreator={\autorFullNameOne},
    pdfsubject={\projektKurztitel -- \projektLangtitel},
    pdfkeywords={\projektKurztitel -- \projektLangtitel},
}

% Einbinden von Bilder (.jpg || .png) ermöglichen
\usepackage{graphicx}

\usepackage[automake, acronym, toc, nonumberlist, nopostdot, nogroupskip, style=super]{glossaries}
\makeglossaries

%Fußzeile und Kopfzeile 
\usepackage{chngcntr} % fortlaufendes Durchnummerieren der Fußnoten

%Literaturverzeichnis
\usepackage[babel,german=quotes]{csquotes} %Empfohlen
\usepackage[backend=biber,
    style=authoryear,isbn=false
]{biblatex}

%Löscht den Inhalt des Feldes 'note'
\AtEveryBibitem{%
    \clearfield{note}%
}

%URL-Umbrüche anpassen
\setcounter{biburllcpenalty}{7000}
\setcounter{biburlucpenalty}{8000}

\addbibresource{04_Literaturverzeichnis/Literaturverzeichnis.bib}

% Für die Darstellung eines Euro-Zeichens.
\usepackage{eurosym}

% Für das Titelblatt 
\usepackage[pages=some]{background}

%Einbinden von einer PDF-Datei
\usepackage{pdfpages}

%Sonstiges
\usepackage{xspace} % sorgt dafür, dass Leerzeichen hinter parameterlosen Makros nicht als Makroendezeichen interpretiert werden

\usepackage[normalem]{ulem} % unterstreichen [\uline], durchstreichen [\sout] etc. % verwendete Packages
%% Möglichkeit der Erweiterung
\newcommand{\addAbbildungQuelle}[5] % Parameter: scale, caption, label, file, Quelle
{

    \begin{minipage}{0.9\textwidth}
        \centering

        \includegraphics[scale=#1]{02_Datein/02_Abbildungen/#4}

        \captionof{figure}[#2]{#2; Quelle: #5}
        \label{#3}

    \end{minipage}

}

\newcommand{\addTabelleScaleTextQuelle}[5] % Parameter: scale (z.B. \columnwidth), caption, label, file, Quelle
{
    \begin{table}[H]
    \centering
    \captionof{table}[#2]{#2; Quelle: #5}
    \label{#3}
    \begin{#1}
    \input{02_Datein/03_Tabellen/#4}
    \end{#1}

    \end{table}
}

% Weiteres

\newcommand{\Vgl}{Vgl.\xspace}

\newcommand{\Anfuehrungsstriche}[1] %inhalt
{\glqq #1\grqq\xspace}

\newcommand{\ZumBeispiel}{z.B.\xspace}

\newcommand{\Gendern}[2]{#1:#2\xspace} %zur Vereinfachung von Prozessen
%Seitenränder
\setlength{\topskip}{\ht\strutbox} % behebt Warnung von geometry
\geometry{a4paper,left=25mm,right=20mm,top=25mm,bottom=35mm}

\usepackage[
	automark, % Kapitelangaben in Kopfzeile automatisch erstellen
	headsepline, % Trennlinie unter Kopfzeile
	ilines % Trennlinie linksbündig ausrichten
]{scrlayer-scrpage}


% Kopfzeile 	---------------------------------------------------------------------
\pagestyle{scrheadings}
\clearpairofpagestyles

%Linker Teil der Kopfzeile
\renewcommand{\headfont}{\normalfont} % Schrift der Kopfzeile
\ihead{\large{\projektKurztitel}\\
	\begin{minipage}{0.93\textwidth}
		\small{\projektLangtitel} \\
	\end{minipage} \\
	\headmark}

%Rechter Teil der Kopfzeile
\chead{}
\ohead{\includegraphics[scale=0.12]{\betriebHSHarzShort}}
\setlength{\headheight}{16.2mm} % Höhe der Kopfzeile

% Fußzeile----------------------------------------------------------------------
\ifoot{\autorFullNameOne}
\cfoot{}
\ofoot{\pagemark}

% Schusterjungen und Hurenkinder vermeiden
\clubpenalty = 10000
\widowpenalty = 10000
\displaywidowpenalty = 10000

% Vermeidet das Fußnoten aufgeteilt werden!
\interfootnotelinepenalty=10000


% Abstand zwischen Nummerierung und Überschrift definieren
\newcommand{\colTableLightBlueSpace}{1.5cm}

% Abschnittsüberschriften im selben Stil wie beim Inhaltsverzeichnis einrücken
\renewcommand*{\othersectionlevelsformat}[3]{
	\makebox[\colTableLightBlueSpace][l]{#3\autodot}
}

%Einrückung des Inhaltsverzeichnisses
\RedeclareSectionCommand[tocindent=0em,tocnumwidth=1em]{section}
\RedeclareSectionCommand[tocindent=1em,tocnumwidth=2em]{subsection}
\RedeclareSectionCommand[tocindent=3em,tocnumwidth=3em]{subsubsection}


%Aufzählungszeichen definieren
\renewcommand{\labelitemi}{\small$\blacktriangleright$}
\renewcommand{\labelitemii}{\small$\blacktriangleright$} % Definitionen zum Aussehen der Seiten
%Farben für Links
\usepackage{xcolor}

\definecolor{MainColor}{RGB}{0,107,179}

\definecolor{colTableLightBlue}{RGB}{0,107,179} 


% Möglichkeit der Erweiterung

%Dummy-Text (Befehl \lipsum[<par-range>]; \lipsum[<par-range>][<sentence-range>])
\usepackage{lipsum}
% Möglichkeit der Erweiterung

\definecolor{colorBeispiel}{RGB}{255,255,255} %Weiß



\begin{document}

%Deckblatt einbinden ------------------------------------------------------
\phantomsection
\pdfbookmark[1]{Deckblatt}{deckblatt}
\begin{titlepage}


    %%Titelseite konfigurieren
    \backgroundsetup{
        scale=1,
        color=black,
        opacity=1,
        angle=0,
        contents={%
                \includegraphics[width=\paperwidth,height=\paperheight]{02_Datein/01_Bilder/00_Base/background_hsHarz.png}
            }%
    }

    %Titelseite anzeigen
    \BgThispage

    \vspace*{16ex} %Mit Exposé-Überschrift 16, sonst 20
    \begin{center}
        \huge{\textbf{\kategorie}}\\
        \vspace*{6ex}

        \huge{\textbf{\projektKurztitel}}\\
        \huge{\projektLangtitel}\\
        \vspace*{6ex}

        \LARGE
        {
            angefertigt an der Hochschule Harz \\ Fachbereich Verwaltungswissenschaften
            \\~\\
            Studiengang: \studiengang
        }

    \end{center}

    \vfill

    %Abstand der Tabelleneinträge 
    \renewcommand{\arraystretch}{1.0}
    \begin{table}[h]
        %Schriftgröße anpassen
        \begin{large}
            %Spaltenbreite anpassen, damit das centering passt
            \begin{tabular}{p{0.5\textwidth}p{0.44\textwidth}}
                \textbf{vorgelegt von:}               & \textbf{angefertigt bei:} \\
                                                      &                           \\
                \autorFullNameOne                     & Erstbetreuer              \\
                \autorAnschriftOne                    & \ansprechpartnerOne       \\
                \autorOrtOne                          &                           \\
                %\autorHandyPrivatOne & \\
                %\autorMailPrivatOne & \\
                                                      &                           \\
                Semester: \autorSemesterNummerOne     & Zweitbetreuer             \\
                Matrikel-Nr.: \autorMatrikelnummerOne & \ansprechpartnerTwo       \\
                E-Mail: \autorMailWorkOne             &                           \\
                                                      &                           \\
                                                      &                           \\
                \textbf{Eingereicht am:} \eingereichtAmDatum
            \end{tabular}
        \end{large}
    \end{table}



\end{titlepage}
\cleardoublepage

%Inhaltsverzeichnis --------------------------------------------------------
\phantomsection
\pagenumbering{Roman}
\pdfbookmark[1]{Inhaltsverzeichnis}{inhalt}
\tableofcontents
\cleardoublepage

%Abkürzungsverzeichnis ----------------------------------------------------- 
% Es werden nur die Abkürzungen aufgelistet, die im Text verwendet werden

\newacronym{itsys}{IT-SYS}{IT-Systemhaus der Bundesagentur für Arbeit}
\setlength{\glsdescwidth}{0.7\textwidth} % Spalte 2 in der Breite anpassen
\renewcommand{\glsnamefont}[1]{\textbf{\textcolor{MainColor}{#1}}} %Abkürzungen fett und farbig.
\renewcommand*{\arraystretch}{1.4}% Abstand zwischen den Einträgen
\printglossary[type=\acronymtype,title={Abkürzungsverzeichnis}]
\clearpage

%Abbildungsverzeichnis -----------------------------------------------------
\phantomsection
\listoffigures
\cleardoublepage

%Tabellenverzeichnis -------------------------------------------------------
\phantomsection
\listoftables
\cleardoublepage


%Inhalt ---------------------------------------------------------------------
\pagenumbering{arabic}
\onehalfspacing %Anpassung an Rahmenbedingungen
\section{Einleitung}
\label{sec:Einleitung}

\subsection{Kapitel in Einleitung}
\label{subsec:Kapitel_in_Einleitung}

Hier stehen Informationen über das \acrshort{itsys}.\footnote{\Vgl \cite[1]{01Kategorie_QUELLE_DasITSystemhausderBundesagenturfurArbeit}.}

Hier ist eine Abbildungen:

\addAbbildungQuelle{0.5}{Vereinfachtes Organigramm der Bundesagentur für Arbeit (Zentrale)}{abb:Vereinfachte_Organigramm_BA}{Vereinfachte_Organigramm_BA.png}{In Anlehnung an \cite{02Kategorie_QUELLE_OrganigrammderZentrale}}
\subsubsection{Unterkapitel in Einleitung}
\label{subsec:Unterkapitel_in_Einleitung}

\addTabelleScaleTextQuelle{normalsize}{Vereinfachte Zeitplanung}{tab:VereinfachteZeitplanung}{VereinfachteZeitplanung.tex}{Eigene Darstellung}

\autoref{tab:VereinfachteZeitplanung} zeigt die vereinfachte Zeitplanung.


\clearpage

%Literatur ------------------------------------------------------------------

\clearpage
\setlength{\bibitemsep}{6pt}
\renewcommand{\refname}{Literaturverzeichnis} %Überschrift abändern
\printbibliography

%Eidesstattliche Erklärung --------------------------------------------------
\input{00_Base/04_EidesstattlicheErklärung_single.tex}
\clearpage


%Anhang ---------------------------------------------------------------------
\clearpage
\appendix
\pagenumbering{roman}
\section{Anhang}
\subsection{Kapitel im Anhang}

\subsubsection{Unterkapitel im Anhang}

Hier könnte Ihre Werbung stehen.

\addAbbildungQuelle{0.5}{Vereinfachtes Organigramm der Bundesagentur für Arbeit (Zentrale) im Anhang}{abb:Vereinfachte_Organigramm_BA}{02_Anhang/Anhang_Vereinfachte_Organigramm_BA.png}{In Anlehnung an \cite{02Kategorie_QUELLE_OrganigrammderZentrale}}

\addTabelleScaleTextQuelle{normalsize}{Vereinfachte Zeitplanung im Anhang}{tab:VereinfachteZeitplanung}{02_Anhang/Anhang_VereinfachteZeitplanung.tex}{Eigene Darstellung}



\end{document}

